\section{Proof-carrying observation graphs}
\label{sec:pcog}

This section defines the central architecture.  The graph is a meta-level
finite object.  Its nodes compile to fixed primitive-recursive symbols and
fixed first-order formulas.  The graph itself is not encoded as a number to be
interpreted by iPRA.

\subsection{Proof-producing rules}

Fix an arity \(n\) for the free numerical parameters and an ambient context
\(\Gamma\) of formulas.  A proof-producing rule has a finite list of required
formulas \(A_1,\ldots,A_r\), an output formula \(B\), and a host-level
constructor
\[
  (\Gamma\vdash A_1)\times\cdots\times(\Gamma\vdash A_r)
  \longrightarrow (\Gamma\vdash B),
\]
where every judgment means \(\Prf\;\PRAAx\).  The constructor may use only the
repository's proof rules and previously proved derived rules.

\begin{definition}[Proof-carrying observation graph]
A PCOG is a finite topologically ordered list of nodes.  Each node specifies:
\begin{enumerate}
\item a finite list of predecessor outputs;
\item concrete primitive-recursive symbols and formulas for its finite data;
\item a local proof-producing rule;
\item optional demand transformers and resource bounds;
\item one exported formula.
\end{enumerate}
The final graph output is one designated node formula.
\end{definition}

The word ``observation'' emphasizes that a node need not construct a canonical
mathematical object.  It constructs enough finite evidence to support later
observations.

\subsection{Five node classes}

The framework normalizes applications into five node classes.

\paragraph{Exact equation nodes.}
An exact node computes a PR term and exports defining equations or a
quantifier-free property.  Examples include tuple projections, polynomial
evaluation, matrix multiplication at fixed dimension, rational normalization,
finite substitution, and restriction of a table.  Its proof uses defining
equations, congruence, and finite algebraic lemmas.

\paragraph{Trace nodes.}
A trace node contains an initializer, step function, quantifier-free invariant,
base derivation, step derivation, and fuel.  It exports the terminal invariant
or a theorem-specific consequence.  Bounded search, bisection, Euclidean
algorithms, coefficient recurrences, dynamic programming, finite products,
and descent procedures all belong here.

\paragraph{Relational cloud nodes.}
A cloud node exports a finite bound and a validity predicate rather than one
selected witness.  It may also export compatibility edges to another cloud.
Backward viability and bounded least search can construct a finite path when
needed, but the cloud remains available for witness-independent stability
arguments.

\paragraph{Observation nodes.}
An observation node compresses certificates or clouds to a quantity that is
stable under the relevant equivalence: an interval, ordered multiset, norm,
residual bound, finite homology table, spectral block, or truth value of a
bounded predicate.  It carries a theorem comparing observations from any two
valid inputs.

\paragraph{Finalizer nodes.}
A finalizer converts terminal finite evidence into the public theorem formula.
Typical examples instantiate a universal trace theorem at a fuel, eliminate
certificate implications, apply a gap lemma, or filter a complete bounded
enumeration.

These node classes are not disjoint implementation types; they are a proof
normal form.  A sophisticated operation can be decomposed into several nodes
so that its hard mathematical estimate is isolated from routine execution.

\subsection{Trace packages}

The central recursive interface is deliberately identical to the checked
repository theorem.

\begin{definition}[Trace package]
For \(k\) parameters and context \(\Gamma\), a trace package is a tuple
\[
  T=(f,g,I,h_I,h_0,h_S),
\]
where:
\begin{itemize}
\item \(f\) is a \(k\)-ary PR initializer;
\item \(g\) is a \((k+2)\)-ary PR step;
\item \(I(s,t,\vec x)\) is a formula with a proof \(h_I\) that it is
quantifier-free;
\item \(h_0\) is an embedded proof of the initialized invariant;
\item \(h_S\) is an embedded proof that one step preserves the invariant.
\end{itemize}
\end{definition}

A trace exit adds a fuel term \(F(\vec x)\), an output formula \(B(\vec x)\),
and an embedded proof
\[
  I(\Run(F,\vec x),F,\vec x),\Gamma\ \vdash\ B(\vec x).
\]
The checked trace theorem and one cut produce \(\Gamma\vdash B\).

This packaging has an important engineering advantage: the generic compiler
never needs to inspect the internal syntax of \(I\)'s proof.  It only verifies
that the required fields have the actual \texttt{Prf} types.

\subsection{Static graph layer and recursive component layer}

A large formal development should not force its entire computation into one
giant induction invariant.  Instead, collapse each genuinely recursive
component to one trace node.  The resulting component graph is acyclic and is
composed by cut.

\begin{principle}[One induction per recursive component]
Use quantifier-free induction exactly where the mathematical algorithm has a
uniform recurrence.  Use topological proof composition everywhere else.
\end{principle}

This gives smaller invariants, better reuse, and cleaner error localization.
For example, a root-factorization proof can have separate trace nodes for
coefficient evaluation, root isolation, deflation, and residual checking.
Their precision demands are connected by the static graph.  A monolithic
invariant over all phases is possible but harder to formalize and audit.

\subsection{Local contract shape}

A deterministic node \(v\) with predecessors \(u_1,\ldots,u_r\) and output
demand \(q\) has a local theorem of the schematic form
\begin{align*}
 &\bigwedge_{j=1}^{r}
 \Valid_{u_j}(x_j,\delta_{v,j}(q,\vec z))
 \ \wedge\ \Cert_v(\vec x,q,c)
 \\
 &\hspace{8em}\longrightarrow
 \Valid_v(E_v(\vec x,q,c),q).
\end{align*}
Here \(\vec z\) denotes parameters and any finite summaries explicitly
allowed to influence demand.  The graph compiler never assumes this formula
because it is true in \(\N\); the node package contains its embedded proof.

For a relational node, the conclusion is
\[
  c<B_v(\vec x,q)\ \wedge\ R_v(\vec x,q,c),
\]
or a finite list containing all such \(c\).  Downstream nodes either consume
all candidates by a bounded fold or request a viable finite path.

For an observation node, the local theorem compares two valid certificates:
\[
 \Cert_v(c_1)\wedge\Cert_v(c_2)
 \longrightarrow d(\Obs(c_1),\Obs(c_2))\leq\varepsilon(q).
\]
This statement deliberately does not require \(c_1=c_2\).

\subsection{Contexts, assumptions, and parameters}

The graph may be compiled under an ambient context \(\Gamma\).  This is
essential for conditional theorems: positivity assumptions, a supplied height
bound, a Jacobian inverse certificate, a cover, or a finite input-validity
condition remain hypotheses.  The final public theorem can then universally
close the numerical parameters and introduce the assumptions as implications.

The graph compiler uses the existing substitution and weakening theorems to
instantiate local rules.  No separate structural metatheory is needed.  The
provided Lean core represents a compiled prefix by a list of output formulas
and a function returning a proof of every member.  Extending the prefix with a
node is a single case split on list membership followed by the node's
proof-producing rule.

\subsection{Why this is more general than a theorem-specific script}

A theorem-specific formalization hard-codes one sequence of computations and
proofs.  A PCOG factors the work along stable interfaces:

\begin{itemize}
\item Boolean folds and bounded search are shared by every finite verifier;
\item trace reflection is shared by every uniform recurrence;
\item backward demand propagation is shared by quantitative constructions;
\item relational viability is shared by non-unique finite approximations;
\item stable-observation comparison is shared across witnesses and producers;
\item height/list conversion is shared by every bounded Diophantine search;
\item graph composition is independent of the mathematics inside nodes.
\end{itemize}

The framework is therefore not a single universal mathematical theorem.  It
is a small collection of compiler theorems whose premises are exactly the
local mathematical obligations that vary by subject.

\subsection{No hidden semantic bridge}

At no point is there a function of type
\[
  (\N\models\varphi)\longrightarrow \Prf\;\PRAAx\;[]\;\varphi.
\]
Such a bridge would be unsound as a general principle.  The only inputs capable
of producing a new object theorem are existing \texttt{Prf} terms, defining
equations furnished by \(\PRAAx\), and applications of its induction axiom.
Python tests, host evaluations, and mathematical prose help design and audit
these inputs; they do not inhabit them.
